% Programmes manquants du E.32 — généré depuis la base de la plateforme
% Compilation : pdflatex note_fonctionnalites.tex (deux passes)
\documentclass[11pt,a4paper]{article}
\usepackage[utf8]{inputenc}
\usepackage[T1]{fontenc}
\IfFileExists{french.ldf}{\usepackage[french]{babel}}{\usepackage{babel}} % babel-french si disponible
\usepackage{lmodern}
\usepackage[margin=2.3cm]{geometry}
\usepackage{xcolor}
\usepackage{booktabs}
\usepackage{array}
\usepackage{longtable}
\usepackage{enumitem}
\usepackage{tcolorbox}
\usepackage{titlesec}
\usepackage{fancyhdr}
\usepackage{hyperref}

\definecolor{nuit}{HTML}{1B2A4A}
\definecolor{or}{HTML}{B8922E}
\definecolor{vert}{HTML}{2E7D5B}
\definecolor{ivoire}{HTML}{F7F4EC}
\definecolor{encre2}{HTML}{5E6677}

\hypersetup{colorlinks=true, linkcolor=nuit, urlcolor=nuit, pdftitle={Plateforme EXETAT — programmes manquants}}
\titleformat{\section}{\color{nuit}\large\bfseries}{\thesection.}{0.6em}{}
\titleformat{\subsection}{\color{nuit}\normalsize\bfseries}{\thesubsection.}{0.6em}{}
\setlist[itemize]{leftmargin=1.4em, itemsep=2pt, topsep=3pt}
\setlength{\parskip}{4pt}
\setlength{\parindent}{0pt}

\pagestyle{fancy}
\fancyhf{}
\fancyhead[L]{\small\color{encre2}Plateforme EXETAT — programmes manquants}
\fancyhead[R]{\small\color{encre2}Novembre 2026}
\fancyfoot[C]{\small\color{encre2}\thepage}
\renewcommand{\headrulewidth}{0.3pt}

\newtcolorbox{encadre}[1][]{colback=ivoire, colframe=or, boxrule=0.6pt, arc=2pt, left=8pt, right=8pt, top=6pt, bottom=6pt, fonttitle=\bfseries\color{nuit}, #1}


\begin{document}
{\color{nuit}\LARGE\bfseries Examen d'État (E.32)}\\[2pt]
{\color{nuit}\Large Matières et options dont les programmes manquent}\\[4pt]
{\color{encre2}État de la base de la plateforme au 25 novembre 2026 (après chargement de la matière Actualité) --- classe terminale des humanités (4\textsuperscript{e} humanités, ex-6\textsuperscript{e} secondaire)}

\vspace{8pt}
\begin{encadre}[title=Vue d'ensemble]
\begin{itemize}
\item \textbf{79 disciplines} du E.32 sans programme chargé, sur 234. La matière « Actualité » est désormais couverte par un programme de quinze thèmes stables (vie nationale, Afrique, monde), la période de référence étant fixée par le jury à chaque session.
\item \textbf{37 épreuves sur 111} non productibles, toutes bloquées par au moins une matière sans programme ; 74 épreuves sont actives.
\item \textbf{Options} (51) : 19 entièrement couvertes, 32 bloquées par une ou plusieurs matières manquantes.
\end{itemize}
\end{encadre}

\section{Options entièrement couvertes (19)}
Agriculture générale, Construction, Coupe et couture, Hydro-pneumatique, Hébergement, Hôtellerie et restauration, Hôtesse d'accueil, Latin-Mathématique, Latin-Philosophie, Menuiserie, Mécanique automobile, Normale, Nutrition, Plomberie (installation sanitaire), Pédagogie générale, Sciences, Sociale, Vétérinaire, Électricité.

\section{Options et matières manquantes}
Pour chaque option : l'épreuve concernée et, en face, les matières dont le programme manque.

\subsection{Humanités littéraires et scientifiques}
\begin{longtable}{@{}>{\raggedright\arraybackslash}p{3.4cm}>{\raggedright\arraybackslash}p{4.4cm}>{\raggedright\arraybackslash}p{7.4cm}@{}}
\toprule
\textbf{Option} & \textbf{Épreuve} & \textbf{Matières sans programme}\\
\midrule\endhead
Latin-Grec & Grec et latin (M.48) & Grec\\
\bottomrule
\end{longtable}

\subsection{Humanités pédagogiques}
\begin{longtable}{@{}>{\raggedright\arraybackslash}p{3.4cm}>{\raggedright\arraybackslash}p{4.4cm}>{\raggedright\arraybackslash}p{7.4cm}@{}}
\toprule
\textbf{Option} & \textbf{Épreuve} & \textbf{Matières sans programme}\\
\midrule\endhead
Pédagogie maternelle & Exploration du milieu et sciences (maternelle) (M.92) & Exploration du milieu\\
Pédagogie pré-scolaire & Exploration du milieu et sciences (maternelle) (M.92) & Exploration du milieu\\
Éducation physique et sportive & Sciences de l'éducation (EPS) (M.50) & Théorie de l'éducation physique et sportive\\
\bottomrule
\end{longtable}

\subsection{Humanités commerciales}
\begin{longtable}{@{}>{\raggedright\arraybackslash}p{3.4cm}>{\raggedright\arraybackslash}p{4.4cm}>{\raggedright\arraybackslash}p{7.4cm}@{}}
\toprule
\textbf{Option} & \textbf{Épreuve} & \textbf{Matières sans programme}\\
\midrule\endhead
Commerciale et gestion & Langues et correspondance commerciale (M.38) & Correspondance commerciale anglaise\\
Secrétariat-Administration & Langues et correspondance commerciale (M.38) & Correspondance commerciale anglaise\\
\bottomrule
\end{longtable}

\subsection{Hôtellerie, accueil et tourisme}
\begin{longtable}{@{}>{\raggedright\arraybackslash}p{3.4cm}>{\raggedright\arraybackslash}p{4.4cm}>{\raggedright\arraybackslash}p{7.4cm}@{}}
\toprule
\textbf{Option} & \textbf{Épreuve} & \textbf{Matières sans programme}\\
\midrule\endhead
Tourisme & Mathématiques et tourisme (M.98) & Photographie\\
\bottomrule
\end{longtable}

\subsection{Humanités agricoles}
\begin{longtable}{@{}>{\raggedright\arraybackslash}p{3.4cm}>{\raggedright\arraybackslash}p{4.4cm}>{\raggedright\arraybackslash}p{7.4cm}@{}}
\toprule
\textbf{Option} & \textbf{Épreuve} & \textbf{Matières sans programme}\\
\midrule\endhead
Foresterie & Foresterie (M.69) & Dendrométrie, Inventaire forestier, Pathologie forestière\\
Industries agricoles & Industries agricoles (M.67) & Toxicologie\\
Pêche et navigation & Pêche et navigation (M.65) & Conservation, Ichtyologie, Ichtyopathologie, Navigation, Pisciculture, Pêche, Technologie et génie\\
\bottomrule
\end{longtable}

\subsection{Humanités artistiques}
\begin{longtable}{@{}>{\raggedright\arraybackslash}p{3.4cm}>{\raggedright\arraybackslash}p{4.4cm}>{\raggedright\arraybackslash}p{7.4cm}@{}}
\toprule
\textbf{Option} & \textbf{Épreuve} & \textbf{Matières sans programme}\\
\midrule\endhead
Arts dramatiques & Arts dramatiques (M.58) & Esthétique théâtrale, Histoire du décor et du costume, Histoire du théâtre\\
Arts plastiques & Arts plastiques (M.57) & Anatomie artistique, Filmologie, Musicologie\\
Coiffure & Culture générale (coiffure) (M.030) & Histoire de la coiffure\\
 & Coiffure (M.902) & Coloration / décoloration, Coupe cheveux, Soins de beauté, Soutien de coiffure, Éducation plastique\\
Esthétique et coiffure & Esthétique et coiffure (M.901) & Cosmétologie, Technologie de la coiffure, Technologie des soins esthétiques, Éducation plastique, Éthique et déontologie\\
Musique & Musique (M.59) & Auditions musicales, Histoire de la musique, Théorie musicale, Transposition\\
\bottomrule
\end{longtable}

\subsection{Humanités techniques industrielles}
\begin{longtable}{@{}>{\raggedright\arraybackslash}p{3.4cm}>{\raggedright\arraybackslash}p{4.4cm}>{\raggedright\arraybackslash}p{7.4cm}@{}}
\toprule
\textbf{Option} & \textbf{Épreuve} & \textbf{Matières sans programme}\\
\midrule\endhead
Aviation civile & Aviation civile (M.83) & Aérotechnique, Navigation aérienne, Opération aérienne, Règlement du transport aérien, Réglementation de la circulation aérienne, Service d'information aéronautique, Sécurité et lutte contre l'incendie\\
 & Mathématiques et aviation (M.31) & Notions de droit aérien\\
Chimie industrielle & Chimie industrielle (M.771) & Chimie industrielle\\
Commutation & Commutation (M.80) & Communication télégraphique, Commutation téléphonique\\
 & Mathématiques et technique (commutation) (M.28) & Transmission télégraphique et numérique, Téléphonie générale\\
Construction métallique & Construction métallique (M.88) & Construction métallique, Traçage de tôlerie et chaudronnerie\\
Dessin de bâtiment & Dessin de bâtiment (M.762) & Histoire d'architecture\\
Géologie & Géologie (M.775) & Analyse des données géologiques\\
 & Mathématiques et technique (géologie) (M.255) & Instruments et équipements de diagraphie de SLE\\
Imprimerie & Imprimerie (M.79) & Composition, Machine offset, Machine typo, Reprographie\\
 & Mathématiques et calculation (imprimerie) (M.27) & Calculation\\
Mine et géologie & Mine et géologie (M.772) & Chimie industrielle, Minéralogie\\
 & Mathématiques et technique (mine et géologie) (M.252) & Technique graphique\\
Mines et carrières & Mathématiques et technique (mines et carrières) (M.254) & Automation des mines, Mécanique des roches, Technique hydro-pneumatique\\
Mécanique dessin & Mécanique dessin (M.84) & Technologie hydropneumatique\\
Mécanique générale & Mathématiques et technique (mécanique) (M.22) & Transmission hydraulique et pneumatique\\
Mécanique machines-outils & Mathématiques et technique (mécanique) (M.22) & Transmission hydraulique et pneumatique\\
Métallurgie & Mathématiques et technique (métallurgie) (M.253) & Chimie industrielle, Métallographie\\
Météorologie & Météorologie (M.82) & Climatologie générale, Radioélectricité, Télécommunication météo\\
 & Mathématiques et météorologie (M.30) & Météorologie générale, Météorologie spéciale, Météorologie tropicale\\
Pétrochimie & Pétrochimie (M.86) & Gestion des stocks pétroliers\\
Radio transmission & Radio transmission (M.81) & Antennes et propagation, Courant porteur, Émission et réception\\
 & Mathématiques et technique (radio) (M.29) & Hyperfréquences\\
Électronique industrielle & Électronique industrielle (M.78) & Hyperfréquence\\
\bottomrule
\end{longtable}

\section{Matières manquantes classées par nombre d'options concernées}
Les matières en tête débloquent plusieurs options à la fois.

\begin{longtable}{@{}>{\raggedright\arraybackslash}p{5cm}c>{\raggedright\arraybackslash}p{8.4cm}@{}}
\toprule
\textbf{Matière} & \textbf{Options} & \textbf{Options concernées}\\
\midrule\endhead
Chimie industrielle & 3 & Chimie industrielle, Mine et géologie, Métallurgie\\
Correspondance commerciale anglaise & 2 & Commerciale et gestion, Secrétariat-Administration\\
Exploration du milieu & 2 & Pédagogie maternelle, Pédagogie pré-scolaire\\
Transmission hydraulique et pneumatique & 2 & Mécanique générale, Mécanique machines-outils\\
Éducation plastique & 2 & Coiffure, Esthétique et coiffure\\
Analyse des données géologiques & 1 & Géologie\\
Anatomie artistique & 1 & Arts plastiques\\
Antennes et propagation & 1 & Radio transmission\\
Auditions musicales & 1 & Musique\\
Automation des mines & 1 & Mines et carrières\\
Aérotechnique & 1 & Aviation civile\\
Calculation & 1 & Imprimerie\\
Climatologie générale & 1 & Météorologie\\
Coloration / décoloration & 1 & Coiffure\\
Communication télégraphique & 1 & Commutation\\
Commutation téléphonique & 1 & Commutation\\
Composition & 1 & Imprimerie\\
Conservation & 1 & Pêche et navigation\\
Construction métallique & 1 & Construction métallique\\
Cosmétologie & 1 & Esthétique et coiffure\\
Coupe cheveux & 1 & Coiffure\\
Courant porteur & 1 & Radio transmission\\
Dendrométrie & 1 & Foresterie\\
Esthétique théâtrale & 1 & Arts dramatiques\\
Filmologie & 1 & Arts plastiques\\
Gestion des stocks pétroliers & 1 & Pétrochimie\\
Grec & 1 & Latin-Grec\\
Histoire d'architecture & 1 & Dessin de bâtiment\\
Histoire de la coiffure & 1 & Coiffure\\
Histoire de la musique & 1 & Musique\\
Histoire du décor et du costume & 1 & Arts dramatiques\\
Histoire du théâtre & 1 & Arts dramatiques\\
Hyperfréquence & 1 & Électronique industrielle\\
Hyperfréquences & 1 & Radio transmission\\
Ichtyologie & 1 & Pêche et navigation\\
Ichtyopathologie & 1 & Pêche et navigation\\
Instruments et équipements de diagraphie de SLE & 1 & Géologie\\
Inventaire forestier & 1 & Foresterie\\
Machine offset & 1 & Imprimerie\\
Machine typo & 1 & Imprimerie\\
Minéralogie & 1 & Mine et géologie\\
Musicologie & 1 & Arts plastiques\\
Mécanique des roches & 1 & Mines et carrières\\
Métallographie & 1 & Métallurgie\\
Météorologie générale & 1 & Météorologie\\
Météorologie spéciale & 1 & Météorologie\\
Météorologie tropicale & 1 & Météorologie\\
Navigation & 1 & Pêche et navigation\\
Navigation aérienne & 1 & Aviation civile\\
Notions de droit aérien & 1 & Aviation civile\\
Opération aérienne & 1 & Aviation civile\\
Pathologie forestière & 1 & Foresterie\\
Photographie & 1 & Tourisme\\
Pisciculture & 1 & Pêche et navigation\\
Pêche & 1 & Pêche et navigation\\
Radioélectricité & 1 & Météorologie\\
Reprographie & 1 & Imprimerie\\
Règlement du transport aérien & 1 & Aviation civile\\
Réglementation de la circulation aérienne & 1 & Aviation civile\\
Service d'information aéronautique & 1 & Aviation civile\\
Soins de beauté & 1 & Coiffure\\
Soutien de coiffure & 1 & Coiffure\\
Sécurité et lutte contre l'incendie & 1 & Aviation civile\\
Technique graphique & 1 & Mine et géologie\\
Technique hydro-pneumatique & 1 & Mines et carrières\\
Technologie de la coiffure & 1 & Esthétique et coiffure\\
Technologie des soins esthétiques & 1 & Esthétique et coiffure\\
Technologie et génie & 1 & Pêche et navigation\\
Technologie hydropneumatique & 1 & Mécanique dessin\\
Théorie de l'éducation physique et sportive & 1 & Éducation physique et sportive\\
Théorie musicale & 1 & Musique\\
Toxicologie & 1 & Industries agricoles\\
Transmission télégraphique et numérique & 1 & Commutation\\
Transposition & 1 & Musique\\
Traçage de tôlerie et chaudronnerie & 1 & Construction métallique\\
Télécommunication météo & 1 & Météorologie\\
Téléphonie générale & 1 & Commutation\\
Émission et réception & 1 & Radio transmission\\
Éthique et déontologie & 1 & Esthétique et coiffure\\
\bottomrule
\end{longtable}

\section{Où trouver ces programmes et comment les charger}
\begin{itemize}
\item \textbf{Sources} : DIPROMAD / EPST (programmes nationaux papier non encore mis en ligne) ; écoles techniques de référence par famille : ITA et ISTA (mines, métallurgie, chimie industrielle, mécanique), Institut de la Gombe et INPP (imprimerie, télécommunications), Régie des voies aériennes et METTELSAT (aviation, météorologie), écoles d'arts (Académie des Beaux-Arts, INA), écoles de pêche et de foresterie (ERAIFT, INERA).
\item \textbf{Format d'import} : un fichier CSV à quatre colonnes --- \texttt{discipline, sous\_discipline, chapitre, poids} --- chargé depuis l'atelier (administrateurs). Les chapitres doivent reprendre les intitulés du programme officiel ; le poids (1 par défaut) règle la part de chaque chapitre dans les épreuves.
\item \textbf{Alternative rapide} : un enseignant de l'option dresse la liste des chapitres de l'année terminale dans un tableau à quatre colonnes ; la source est alors marquée « à confirmer » jusqu'à validation par un inspecteur.
\item Dès qu'une matière est chargée, les épreuves qui en dépendent s'activent automatiquement pour la production.
\end{itemize}

\section*{Annexe --- programme de la matière Actualité}
Quinze thèmes stables, en trois sous-disciplines ; la fenêtre temporelle est fixée par le jury à chaque session (par défaut, du 1\textsuperscript{er} janvier de l'année précédente à la date de génération). Les questions ne portent que sur des faits vérifiés par recherche documentaire au moment de la génération, avec source citée dans la correction.

\begin{longtable}{@{}>{\raggedright\arraybackslash}p{3cm}>{\raggedright\arraybackslash}p{11.6cm}c@{}}
\toprule
\textbf{Sous-discipline} & \textbf{Thème} & \textbf{Poids}\\
\midrule\endhead
Vie nationale (RDC) & Institutions et vie politique nationale : Présidence, Gouvernement, Parlement, CENI, Cour constitutionnelle, réformes en cours, calendrier électoral & 4\\
Vie nationale (RDC) & Paix, sécurité et situation dans l'Est : processus de paix, accords, MONUSCO et forces régionales, déplacés & 4\\
Vie nationale (RDC) & Économie et finances publiques : budget, franc congolais, mines et contrats, ZLECAf, inflation, prix des produits de base & 3\\
Vie nationale (RDC) & Société, santé et éducation : gratuité de l'enseignement, EXETAT et réformes, épidémies et campagnes de vaccination, couverture santé universelle & 3\\
Vie nationale (RDC) & Infrastructures, énergie et environnement : routes, barrages dont Inga, forêt du bassin du Congo, COP et crédits carbone, catastrophes naturelles & 2\\
Vie nationale (RDC) & Culture, sport et personnalités congolaises : Léopards, rumba classée UNESCO, prix et distinctions, disparitions marquantes & 2\\
Afrique & Union africaine et organisations régionales : UA, SADC, CEEAC, EAC, CIRGL (sommets, présidences, décisions) & 3\\
Afrique & Vie politique et crises du continent : élections, transitions, coups d'État, conflits, médiations & 3\\
Afrique & Économie et intégration africaines : ZLECAf, monnaies, dette, grands projets, ressources naturelles & 2\\
Afrique & Sport et culture africains : CAN, Jeux africains, prix littéraires et musicaux, cinéma & 2\\
Monde & Organisations internationales : ONU et ses agences, Conseil de sécurité, OMS, FMI et Banque mondiale, Francophonie, CPI (dirigeants, décisions, rapports marquants) & 3\\
Monde & Grandes crises et relations internationales : guerres en cours, négociations de paix, sanctions, élections majeures, grandes puissances et Afrique & 3\\
Monde & Économie mondiale : croissance, pétrole, cours des métaux stratégiques (cobalt, cuivre, lithium), énergie, chaînes d'approvisionnement & 2\\
Monde & Sciences, techniques et environnement : intelligence artificielle, espace, climat et rapports du GIEC, santé mondiale, prix Nobel & 2\\
Monde & Sport et culture mondiaux : Coupe du monde, Jeux olympiques, grands prix culturels, disparitions de personnalités & 2\\
\bottomrule
\end{longtable}
\end{document}